\documentclass{ltc-guji}
\setmainfont{TW-Kai}
\关闭分页

\title{表格测试}
\chapter{世系表}

\begin{document}

% ============================================================
% 测试一：表格嵌入正文（基本三栏表格）
% ============================================================
\begin{正文}

测试一基本表格

\begin{表格}[分栏=3, 栏间距=3pt]

\单元格[行宽=3]{姓名}
\单元格[行宽=3]{世次}
\单元格{备注}

\换栏

\单元格[行宽=3]{黄帝}
\单元格[行宽=3]{第一世}
\单元格{有熊氏}

\换栏

\单元格[行宽=3]{颛顼}
\单元格[行宽=3]{第二世}
\单元格{高阳氏}

\end{表格}

表格后的正文内容

\end{正文}

\clearpage

% ============================================================
% 测试二：两列表格嵌入正文
% ============================================================
\chapter{书目表}

\begin{正文}

古籍书目如下

\begin{表格}[分栏=3, 栏间距=5pt]

\单元格[行宽=4]{书名}
\单元格{作者}

\换栏

\单元格[行宽=4]{史记}
\单元格{司马迁撰}

\换栏

\单元格[行宽=4]{汉书}
\单元格{班固撰}

\end{表格}

以上为书目

\end{正文}

\clearpage

% ============================================================
% 测试三：多栏表格（6行）
% ============================================================
\chapter{朝代表}

\begin{正文}

历代朝代表

\begin{表格}[分栏=6, 栏间距=2pt]

\单元格[行宽=2]{朝代}
\单元格[行宽=2]{都城}
\单元格{年数}

\换栏

\单元格[行宽=2]{夏}
\单元格[行宽=2]{安邑}
\单元格{四百年}

\换栏

\单元格[行宽=2]{商}
\单元格[行宽=2]{殷}
\单元格{六百年}

\换栏

\单元格[行宽=2]{周}
\单元格[行宽=2]{镐京}
\单元格{八百年}

\换栏

\单元格[行宽=2]{秦}
\单元格[行宽=2]{咸阳}
\单元格{十五年}

\换栏

\单元格[行宽=2]{汉}
\单元格[行宽=2]{长安}
\单元格{四百年}

\end{表格}

右表记历代朝代

\end{正文}

\clearpage

% ============================================================
% 测试四：自定义栏高度
% ============================================================
\chapter{栏高度测试}

\begin{正文}

自定义栏高度

\begin{表格}[分栏=3, 栏间距=3pt, 栏高度={3cm, 6cm, 8cm}]

\单元格[行宽=4]{表头}
\单元格{说明}

\换栏

\单元格[行宽=4]{第一行}
\单元格{栏高六厘米}

\换栏

\单元格[行宽=4]{第二行}
\单元格{栏高八厘米}

\end{表格}

栏高度测试完成

\end{正文}

\clearpage

% ============================================================
% 测试五：行数限制
% ============================================================
\chapter{行数测试}

\begin{正文}

行数限制测试

\begin{表格}[分栏=2, 栏间距=3pt, 行数=10]

\单元格[行宽=4]{书名}
\单元格{作者}

\换栏

\单元格[行宽=4]{史记}
\单元格{司马迁}

\end{表格}

行数测试完成

\end{正文}

\clearpage

% ============================================================
% 测试六：栏边框=true，行边框=false（只画栏分隔线，不画行线）
% ============================================================
\chapter{栏边框测试}

\begin{正文}

栏边框开启行边框关闭

\begin{表格}[分栏=3, 栏间距=3pt, 栏边框=true, 行边框=false]

\单元格[行宽=3]{姓名}
\单元格[行宽=3]{世次}
\单元格{备注}

\换栏

\单元格[行宽=3]{黄帝}
\单元格[行宽=3]{第一世}
\单元格{有熊氏}

\换栏

\单元格[行宽=3]{颛顼}
\单元格[行宽=3]{第二世}
\单元格{高阳氏}

\end{表格}

栏边框测试完成

\end{正文}

\clearpage

% ============================================================
% 测试七：栏边框=false，行边框=true（只画行线，不画栏分隔线）
% ============================================================
\chapter{行边框测试}

\begin{正文}

行边框开启栏边框关闭

\begin{表格}[分栏=3, 栏间距=3pt, 行边框=true, 栏边框=false]

\单元格[行宽=3]{朝代}
\单元格[行宽=3]{都城}
\单元格{年数}

\换栏

\单元格[行宽=3]{夏}
\单元格[行宽=3]{安邑}
\单元格{四百年}

\换栏

\单元格[行宽=3]{商}
\单元格[行宽=3]{殷}
\单元格{六百年}

\end{表格}

行边框测试完成

\end{正文}

\clearpage

% ============================================================
% 测试八：两种边框都关闭
% ============================================================
\chapter{无边框测试}

\begin{正文}

行边框和栏边框都关闭

\begin{表格}[分栏=3, 栏间距=3pt, 行边框=false, 栏边框=false]

\单元格[行宽=3]{书名}
\单元格{作者}

\换栏

\单元格[行宽=3]{史记}
\单元格{司马迁}

\换栏

\单元格[行宽=3]{汉书}
\单元格{班固}

\end{表格}

无边框测试完成

\end{正文}

\clearpage

% ============================================================
% 测试九：单元格纵对齐 + 八角墨围反白（textbox 居中）
% ============================================================
\chapter{纵对齐测试}

\begin{正文}

纵对齐测试

\begin{表格}[分栏=3, 栏高度={6cm, 6cm, 6cm}]
\栏格式[行边框=false]
\单元格[行宽=1, 纵对齐=center]{\八角墨围反白{第一世}}
\单元格[行宽=5]{黄帝}

\换栏

\单元格[行宽=1, 纵对齐=center]{\八角墨围反白{第二世}}
\单元格[行宽=5]{颛顼}

\换栏

\单元格[行宽=1, 纵对齐=center]{\八角墨围反白{第三世}}
\单元格[行宽=5]{帝喾}

\end{表格}

纵对齐测试完成

\end{正文}

\clearpage

% ============================================================
% 测试十：Parallel 模式跨页（某栏内容超过一页高度）
% ============================================================
\chapter{跨页表格}

\begin{正文}

跨页表格测试

\begin{表格}[分栏=6, 栏高度={35pt}, 行边框=false, 行填充=页]
\栏格式[行边框=false, 上填充=0pt]
\单元格[行宽=1]{世}
\单元格{姓名与注}

\换栏

\单元格[行宽=1]{一}
\单元格{黄帝}
\单元格{有熊氏之子}
\单元格{黄帝}
\单元格{有熊氏之子}
\单元格{黄帝}
\单元格{有熊氏之子}
\单元格{黄帝}
\单元格{有熊氏之子}
\单元格{黄帝}
\单元格{有熊氏之子}
\单元格{黄帝}
\单元格{有熊氏之子}
\单元格{黄帝}
\单元格{有熊氏之子}
\单元格{黄帝}

\换栏

\单元格[行宽=1]{二}
\单元格{颛顼}
\单元格{高阳氏}

\换栏

\单元格[行宽=1]{三}
\单元格{帝喾}
\单元格{高辛氏}

\换栏

\单元格[行宽=1]{四}
\单元格{唐尧}
\单元格{陶唐氏}

\换栏

\单元格[行宽=1]{五}
\单元格{虞舜}
\单元格{有虞氏}

\end{表格}

跨页表格测试完成

\end{正文}

\end{document}
